% Section 66, visually collated across 134-4/093--100.
\subsection*{66. Asphericity, universal equivalences\newline and the revision of interval terminology (pp. 180--187)}
The entry of 9.4 begins by reassessing the long route to the ``key result'' of section~65. Grothendieck wants to compare a category with a presheaf category and to construct canonical modelizers; the restriction to functors taking values in $(\mathrm{Cat})$ has made that route indirect. Its psychological gain has been the realization that contractible objects and contractibility might recover the weak equivalences. He returns to the program of section~56, specifically part d) on p.~138 after the work on b), while acknowledging that the preceding development has not followed his original priorities (pp.~180--181). The next section will immediately reconsider the emphasis again: aspheric maps organize comparison between modelizers, whereas the contractibility program addresses a particular canonical construction. The terminology on these pages should therefore be read as an active revision, with its dated qualifications, rather than as a settled list of interchangeable definitions. The geometrical uses of asphericity for spaces, topoi and categories are to govern its extension to new settings; the earlier discussion of topoi on p.~37 supplies a reference point. For categories, ordinary asphericity of a functor is tested after the maps $Y/y\to Y$, whereas universal asphericity asks for the corresponding weak equivalence after every base change $Y'\to Y$. These conditions agree for an object viewed over the final category. Their analogues relative to a basic localizer remain subject to the localizer hypotheses in use (pp.~181--182). A presheaf $F$ on $A$ is called $\underline W$-aspheric when its category of elements $A/F$ is aspheric; weak equivalences of presheaves are likewise detected through the categories of elements. The localization axiom implies that an aspheric map in $A^\wedge$ is universally a $\underline W_A$-equivalence with respect to base changes in $A^\wedge$. This quantifier ranges over the presheaf category; it does not automatically assert the result for every arbitrary base change of the corresponding functor in $(\mathrm{Cat})$ (pp.~182--183). Three tests must consequently be kept distinct: absolute asphericity of $A/F\to e$, relative asphericity of $A/F\to A$, and the weaker assertion that this latter functor belongs to $\underline W$. For arbitrary $A$, the first two do not imply one another. With $A$ the discrete category on two objects, a presheaf supported by a singleton at just one object has an absolutely aspheric category of elements, but its projection fails the relative test at the other object. Conversely, the terminal presheaf has identity projection to $A$, hence is relatively aspheric, although its category of elements is the nonaspheric two-point category. If $A$ is aspheric, the first and third tests agree by the two-out-of-three property. The terminal presheaf and the terminal category cannot be interchanged in this argument (pp.~183--184). In a general pair $(M,W)$ Grothendieck now reserves a separate name for the maps universally in $W$, writing $UW\subset W$, and withholds the word aspheric until an appropriate class of localizing maps has been specified. Universal equivalences contain isomorphisms and are stable under composition, but need not themselves be saturated. The handwritten revision from ``is'nt'' to ``need not be'' changes an asserted universal failure into a possibility, and must be applied when reading the final claim (p.~184). The related terminology of contractors is also reconsidered. He proposes reserving contractor for a stronger notion suited to strict test categories, and renaming the earlier contractors of p.~120 multiplicative intervals with a left unit and a left zero. The source's phrase ``left zero object'' is retained. The proposed alternative intersector does not adequately describe an internal-Hom example. These changes of name do not withdraw the contraction argument of section~51 (p.~185). The discussion then returns to the \v Cech condition L e) of p.~171 and reduces its elementary two-piece form to the category $C_o=(b\leftarrow a\to c)$. A cover by two cribles supplies a functor to $C_o$; the asphericity conditions on the cribles and their intersection make that functor aspheric, and the relevant localization axiom makes it a weak equivalence. The source leaves an empty slot between $C'$ and $C''$ in one statement where the comparison witness supplies an intersection sign. The mathematical role is that of an intersection, but the readers preserve the actual blank. The functor $C_o\to\Delta_1$ carrying $a,b$ to $0$ and $c$ to $1$ is fibering, with fibers $\Delta_1$ and $\Delta_o$. The fibration axiom L 5), together with the interval asphericity requirement, therefore gives this Mayer--Vietoris form of L e). The suggestion that L 5) should imply all reasonable \v Cech criteria is an expectation beyond the argument just supplied, not a proved general theorem (pp.~185--186). The connectedness condition previously discussed on p.~165 is next promoted to L 6): $f\in\underline W$ must induce a bijection on $\pi_o$. Equivalently, the component functor factors through $(\mathrm{Hot}_{\underline W})$. Its limited role in the earlier key result no longer seems a reason to treat it as accessory; Grothendieck explicitly distinguishes this requirement from the weaker statement that the associated homotopy structure is the canonical one on $(\mathrm{Cat})$ (p.~186). The entry of 11.4 finally abandons the qualifier weak in ``weak homotopy interval''. A homotopy interval for a specified structure $h$ is henceforth simply an interval whose two endpoint sections are $h$-homotopic. When $h=h_W$, this is broader than the earlier p.~132 requirement that $I\to e$ lie in $UW$. The stronger condition remains available and useful, but it no longer defines the unqualified term. The requirement $\mathrm{Ker}(\delta_o,\delta_1)=\emptyset_M$ is renamed separating rather than separated, to avoid confusion with a separate notion of separation for objects of $M$ (pp.~186--187). The source evidence is consequential here: ``from now on'', ``page 132'' and ``till this very minute'' are typed interlinear additions; their pen placing curves do not cancel the adjoining live words. In the opening line of p.~187 only the n of an is canceled, leaving a live, and the two quotation pairs around homotopy interval have separate machine cancellations. The readers retain doubtful overprint layers, the clipped canceled cell on p.~184, the unresolved oval/parenthesis patch on p.~185, and the recurring graphic marks around h without assigning them an unsupported pen provenance. The isolated typed slips gontext and notes on p.~184 are treated in the corrected forms only, using the same words correctly typed on other reviewed pages; neither those corrections nor the explanatory distinctions above introduce new mathematical hypotheses into the source.
